% r03_hyb_domdiff: proof of \Cref{rot:prop:BD}, part III: the domain-difference lemma (two tori, or a torus against the infinite lattice) and the restricted precision inverses on unions of big blocks.

\begin{lemma}[domain difference]\label{rot:hyb:lem:domdiff}
Let $\Omega_1=\Lambda_X/P_1$ be a torus fitted by both lattices with $s_\infty(P_1)\ge10L$ in unit scale, and $F$ a fundamental region of $P_1$ that is a union of $(r+1)$-blocks. Let $\Omega_2$ be $\Lambda_X$, or a torus $\Lambda_X/P_2$ fitted by both lattices into which $F$ embeds (the projection is injective on $F$); $F$ is identified with its images. Let $A_i$ be regular on $\Omega_i$ with $e_r^\beta\le c_0(a)$, and let $A_1=A_2$ on every bond with both ends in $F_0\defeq\{x\in F:\dist(x,\partial F)\ge W_0\}$, for a fixed $W_0\ge0$. Put $H_i\defeq-\Delta_{A_i}+m_r^2+a_{r+1}L^{-2}P_{r+1}(A_i)$ on $\Omega_i$ (the $(r+1)$-pin, $P_{r+1}=Q_{r+1}^*Q_{r+1}$ with the level-$(r+1)$ blocks) and $G_i\defeq H_i^{-1}$. Let $d_i$ be the quotient metric of $\Omega_i$ and $d_\wedge\defeq\min\{d_1,d_2\}$ on $F$; distances to $\partial F$ and to $F_0^c$ are taken in $\Lambda_X$, for a lift of $F$. There are $C,\delta>0$ depending on $a,L,W_0$ only such that for $f,f'$ supported in unit cubes in $F$
\begin{equation}\label{rot:hyb:eq:B1}
|\ip{f'}{(G_1-G_2)f}|\le C\exp\!\big[-\delta\big(d_\wedge(\supp f,\supp f')+\dist(\supp f,\partial F)+\dist(\supp f',\partial F)\big)\big]\|f\|\|f'\|,
\end{equation}
and, for $y,y'\in\Gamma_r\cap F$,
\begin{equation}\label{rot:hyb:eq:B2}
|C^{(r)}(\Omega_1,A_1;y,y')-C^{(r)}(\Omega_2,A_2;y,y')|\le C\exp\!\big[-\delta\big(d_\wedge(y,y')+\dist(y,\partial F)+\dist(y',\partial F)\big)\big].
\end{equation}
The same bound \eqref{rot:hyb:eq:B2} holds for $\Delta^{(r)}=A^X$, from \eqref{rot:hyb:eq:B1} at the $r$-pin. In particular $\delta C^{(r)}$ and $\delta\Delta^{(r)}$ carry the factor $e^{-\delta\dist(\{y,y'\},\partial F)}$ of \cite[Thm.~3.3]{K1}, in the metric $d_\wedge$.
\end{lemma}

The metric is forced: for $y$ near one face of $F$ and $y'$ near the opposite face, $d_1(y,y')=O(1)$, so $C^{(r)}(\Omega_1,A_1;y,y')$ is of order one, while $C^{(r)}(\Omega_2,A_2;y,y')$ is small when $\Omega_2=\Lambda_X$; no bound of the form \eqref{rot:hyb:eq:B2} with the distance $|y-y'|$ of the lift in place of $d_\wedge$ can hold. The kernel bounds on each domain separately are in the quotient distance of that domain (\Cref{rot:hyb:prop:CT}).

\begin{proof}
\emph{Step 0 (decay on each domain).} By \Cref{rot:hyb:prop:CT} at the $(r+1)$-pin, with the distance weight on a torus and the linear weight on $\Lambda_X$, for $f,f'$ supported in unit cubes and $i=1,2$,
$|\ip f{G_if'}|,\ |\ip f{D_{A_i}G_if'}|,\ |\ip f{G_iD_{A_i}^*f'}|\le c_1e^{-\theta_0d_i(\supp f,\supp f')}\|f\|\|f'\|$, $\theta_0=\theta_0(a,L)$.

\emph{Step 1 (locality).} Put $r_1\defeq1+\sqrt3L$ and $S\defeq\{x\in F_0:\dist(x,F_0^c)\ge r_1\}$. For $u$ supported in $S$, $H_1u=H_2u$ (both supported in $F_0$, identified). Indeed $(H_iu)(x)$ involves $u$ at the bond neighbours of $x$ and on the $(r+1)$-block of $x$, the bond variables of the bonds meeting $\supp u$ (both ends in $F_0$), and the block functions $w_z$ of the blocks meeting $S$; such a block lies in $F$ ($F$ is a union of $(r+1)$-blocks) within $\sqrt3L$ of $S$, hence in $F_0$, where $A_1=A_2$; its contours lie inside it (\Cref{rot:hyb:lem:alg}), so the contour phases of $w_z$ agree.

\emph{Step 2 (geometric resolvent identity).} Let $\chi(x)\defeq\min\{1,\max\{0,\dist(x,F_0^c)-r_1\}\}$ (Lipschitz constant $1$), so $\supp\chi\subset S$ and $\chi=1$ on $S''\defeq\{x:\dist(x,F_0^c)\ge r_1+1\}$. Let $f,f'$ be supported in $S''$. Then $\chi G_2f$ is supported in $S$, so by Step~1 $H_1(\chi G_2f)=H_2(\chi G_2f)=\chi f+[H_2,\chi]G_2f=f+[H_2,\chi]G_2f$. The function $[H_2,\chi]G_2f$ vanishes at points farther than $r_1$ from the set where $\chi$ is not locally constant, so it is supported in $F_0$ and is a function on $\Omega_1$. Applying $G_1$ and pairing with $f'$ ($\chi f'=f'$, $\chi$ real),
\begin{equation}\label{rot:hyb:eq:B3}
\ip{f'}{G_2f}-\ip{f'}{G_1f}=\ip{f'}{G_1[H_2,\chi]G_2f}.
\end{equation}

\emph{Step 3 (the commutator).} With $(D_A\phi)(b)=\eta'^{-1}(U(-A_b)\phi(b_+)-\phi(b_-))$, $D_A(\chi u)(b)=\chi(b_-)(D_Au)(b)+(Mu)(b)$, $(Mu)(b)\defeq\eta'^{-1}(\chi(b_+)-\chi(b_-))U(-A_b)u(b_+)$; taking adjoints, $\chi D_A^*-D_A^*\chi_-=M^*$, hence $[-\Delta_A,\chi]=D_A^*M-M^*D_A$. Since $|\chi(b_+)-\chi(b_-)|\le\eta'$ and every site is the end $b_+$ of three bonds, $\|M\|\le\sqrt3$. For the pin, $P_{r+1}=\sum_zw_zw_z^*$ with $\|w_z\|=1$ and $\supp w_z\subset B_{r+1}(z)$ (\Cref{rot:hyb:lem:alg}), so $[P_{r+1},\chi]=\sum_z\big(w_zw_z^*(\chi-\chi(z))-(\chi-\chi(z))w_zw_z^*\big)$, each term acting in $\ell^2(B_{r+1}(z))$ with norm $\le2\sup_{B_{r+1}(z)}|\chi-\chi(z)|\le\sqrt3L$, and zero unless $\chi$ is non-constant on the block; $m_r^2$ commutes with $\chi$. All these terms live on bonds and blocks meeting the slab $\{r_1\le\dist(\cdot,F_0^c)\le r_1+1\}$, hence on the set $Y\defeq\{x:\tfrac12\le\dist(x,F_0^c)\le r_1+2+\sqrt3L\}\subset F_0$ (the pieces lie within $\sqrt3L+\eta'$ of the slab and $r_1-\sqrt3L-\eta'>\tfrac12$), on whose bonds $D_{A_1}=D_{A_2}$.

\emph{Step 4 (bound of \eqref{rot:hyb:eq:B3}).} Write $\ip{f'}{G_1D^*MG_2f}=\ip{DG_1f'}{MG_2f}$, $\ip{f'}{G_1M^*DG_2f}=\ip{MG_1f'}{DG_2f}$ and the pin terms as $\ip{G_1f'}{[w_zw_z^*,\chi]G_2f}$. Split $Y$ into unit cubes $Q$ (every bond and block meets at most $(\sqrt3L+2)^3$ of them) and apply the Cauchy--Schwarz inequality cube by cube, with $\|\mathbf 1_QDG_2f\|=\sup\{|\ip h{DG_2f}|:h\in\ell^2(Q),\|h\|=1\}$ and the same without $D$, and Step~0 on $\Omega_1$ for the left factor and on $\Omega_2$ for the right factor. Since $\dist(x,\partial F)\le\dist(x,F_0^c)+W_0$, $Y\subset K\defeq\{x\in F:\dist(x,\partial F)\le W_0+r_2\}$ with $r_2\defeq2r_1+2$, and
\begin{equation}\label{rot:hyb:eq:B4}
|\ip{f'}{(G_2-G_1)f}|\le C_L\sum_{Q\subset K}e^{-\theta_0(d_1(\supp f',Q)+d_2(Q,\supp f))}\|f\|\|f'\|.
\end{equation}
\emph{Step 5 (geometry).} For $x\in F$ and $q\in F$, a shortest lattice path in $\Omega_i$ from $x$ to $q$, lifted from $x$, either stays in the lift of $F$ (the projection is injective on $F$), and then has length $\ge|x-q|\ge\dist(x,\partial F)-\dist(q,\partial F)$, or it leaves it, and then has length $\ge\dist(x,\partial F)$. Hence $d_i(x,q)\ge\dist(x,\partial F)-\dist(q,\partial F)$, and for $q\in K$, $d_i(x,q)\ge\dist(x,\partial F)-W_0-r_2-1$ (the $1$ for the unit cubes). Splitting $\theta_0(d_1+d_2)=\frac{\theta_0}2d_1+\frac{\theta_0}2(d_1+d_2)$ in \eqref{rot:hyb:eq:B4},
\[
\theta_0\big(d_1(\supp f',Q)+d_2(Q,\supp f)\big)\ge\tfrac{\theta_0}2d_1(\supp f',Q)+\tfrac{\theta_0}2\big(\dist(\supp f,\partial F)+\dist(\supp f',\partial F)\big)-\theta_0(W_0+r_2+3),
\]
and $\sum_Qe^{-(\theta_0/2)d_1(\supp f',Q)}\le C(\theta_0)$, since the unit cubes of a fundamental region at $d_1$-distance in $[n,n+1)$ from a unit cube number at most $C(n+2)^3$ (choose for each the lift nearest to a fixed lift of $\supp f'$). So, for $f,f'$ supported in $S''$,
\[
\text{(ii)}\qquad|\ip{f'}{(G_1-G_2)f}|\le C'e^{-(\theta_0/2)(\dist(\supp f,\partial F)+\dist(\supp f',\partial F))}\|f\|\|f'\|.
\]
Since $\dist(x,F_0^c)\ge\dist(x,\partial F)-W_0$, a unit cube not contained in $S''$ lies within $r_3\defeq W_0+r_1+3$ of $\partial F$. If $\supp f$ is such a cube, then by the first inequality of this step $d_\wedge(\supp f,\supp f')\ge\dist(\supp f',\partial F)-r_3-2$, so $\frac{\theta_0}2(\dist(\supp f,\partial F)+\dist(\supp f',\partial F))\le\frac{\theta_0}2d_\wedge(\supp f,\supp f')+\theta_0(r_3+1)$, and (ii) follows from (i) below with $C'$ replaced by $\max\{C',2c_1e^{\theta_0(r_3+1)}\}$; the same if $\supp f'$ is such a cube. So (ii) holds for all $f,f'$ supported in unit cubes in $F$.

\emph{Step 6 (proof of \eqref{rot:hyb:eq:B1}).} By Step~0, (i) $|\ip{f'}{(G_1-G_2)f}|\le|\ip{f'}{G_1f}|+|\ip{f'}{G_2f}|\le2c_1e^{-\theta_0d_\wedge(\supp f,\supp f')}\|f\|\|f'\|$. Then $\min\{\text{(i)},\text{(ii)}\}\le\text{(i)}^{1/2}\text{(ii)}^{1/2}$, whose exponent is $\frac{\theta_0}2d_\wedge+\frac{\theta_0}4(\dist(\supp f,\partial F)+\dist(\supp f',\partial F))$; this is \eqref{rot:hyb:eq:B1} with $\delta=\theta_0/4$ and $C=(2c_1C')^{1/2}$.

\emph{Step 7 (proof of \eqref{rot:hyb:eq:B2}).} Take $\Lambda=\Omega^{(r)}$ in \cite[(3.5)]{Ba4}; then $B^r(\Lambda)=\Omega$ and $G_r(\Omega_i,\Omega_i^{(r)},A_i)=G_i$. In the conventions above ($Q_rQ_r^*=\mathbf 1$, $QQ^*=\mathbf 1$, $P=Q^*Q$, $P_{r+1}=Q_r^*PQ_r$, with $Q_r=Q_X(A)$ and $Q=Q^{(1)}(A)$), \cite[(3.6)]{Ba4} reads
\[
C^{(r)}(\Omega_i,A_i)=V_i^*G_iV_i-\tfrac{\beta}{a_r}P(A_i)+a_r^{-1},\qquad V_i\defeq Q_r^*(A_i)\big(1-\beta P(A_i)\big),\qquad\beta\defeq\frac{a_{r+1}L^{-2}}{a_r},
\]
with $a_{r+1}=aa_r/(aL^{-2}+a_r)$; it is obtained by integrating the coarse variable out of the Gaussian of \cite[(3.1)]{Ba4}, is an algebraic identity between bounded operators, and holds on $\Lambda_X$ as on a torus. The vector $v_i(y)\defeq V_i\delta_y$ is supported in the $(r+1)$-block of $y$, is built from contours inside that block (\Cref{rot:hyb:lem:alg}), and has $\|v_i(y)\|\le1+\beta\le2$.
\emph{Case 1:} the $(r+1)$-blocks of $y$ and $y'$ lie in $F_0$. Then $v_1(y)=v_2(y)=:v(y)$, the same for $y'$, and $P(A_1;y,y')=P(A_2;y,y')$, since $A_1=A_2$ there; so $\delta C^{(r)}(y,y')=\ip{v(y)}{(G_1-G_2)v(y')}$. Splitting $v(y)$, $v(y')$ into at most $(\sqrt3L+2)^3$ unit-cube pieces and applying \eqref{rot:hyb:eq:B1} to each pair, $|\delta C^{(r)}(y,y')|\le4(\sqrt3L+2)^6e^{6\delta\sqrt3L}C\,e^{-\delta(d_\wedge(y,y')+\dist(y,\partial F)+\dist(y',\partial F))}$.
\emph{Case 2:} a $(r+1)$-block of $y$ or $y'$ meets $F\setminus F_0$, say that of $y$, so $\dist(y,\partial F)\le r_4\defeq W_0+\sqrt3L$. Then $|\delta C^{(r)}|\le|C^{(r)}(\Omega_1)|+|C^{(r)}(\Omega_2)|\le2C_0e^{-\delta_0d_\wedge(y,y')}$ by the single-domain decay of \Cref{rot:hyb:cor:Ck} on each domain (in its quotient distance), and $d_\wedge(y,y')\ge\dist(y',\partial F)-r_4$ by Step~5, so $e^{-\delta_0d_\wedge}\le e^{\delta_0r_4}e^{-(\delta_0/2)(d_\wedge(y,y')+\dist(y,\partial F)+\dist(y',\partial F))}$. Both cases give \eqref{rot:hyb:eq:B2} with $\min\{\delta,\delta_0/2\}$ in place of $\delta$.
For $\Delta^{(r)}=a_r-a_r^2Q_rG^X_rQ_r^*$ (\cite[(2.14)]{K1}, the $r$-pin) the same argument applies with $w_z$ in place of $v(y)$, with \eqref{rot:hyb:eq:B1} at the $r$-pin (Steps 0--6 with $r_1=1+\sqrt3$, \Cref{rot:hyb:lem:coer}).
\end{proof}

\begin{proposition}[{\cite[Prop.~5.6]{K2} for the hybrid flow}]\label{rot:hyb:prop:CLam}
Let $X\defeq A^X(A)+aL^{-2}P(A)$ on the coarse lattice $\Gamma_r$ of the domain, $P=Q^{(1)}(A)^*Q^{(1)}(A)$. Then $(L^2\gamma_0^{-1}+a_r^{-1})^{-1}\le X\le a_r+aL^{-2}$ and $|X(y,y')|\le Ce^{-\theta d(y,y')}$; for every subset $\Lambda$ of the coarse lattice, in particular every union of big blocks, the same bounds hold for the compression $X|_\Lambda$, and $\hat C^{(r)}_\Lambda\defeq(X|_\Lambda)^{-1}$ (\cite[(5.53)]{K2}) obeys $|\hat C^{(r)}_\Lambda(y,y')|\le Ce^{-t_*d(y,y')}$ (\cite[(5.54)]{K2}), $|\hat C^{(r)}_\Lambda(y,y')-C^{(r)}(y,y')|\le Ce^{-(t_*/6)(d(y,y')+d(y,\Lambda^c)+d(y',\Lambda^c))}$, and, under the hypotheses of \Cref{rot:hyb:lem:domdiff} with $\Lambda\subset F$, the two-domain bound \eqref{rot:hyb:eq:B2} in the metric $d_\wedge$. The constants depend on $a$, $L$ (and $W_0$) only.
\end{proposition}
\begin{proof}
Upper bound: $A^X=(a_r^{-1}+Q_rC_AQ_r^*)^{-1}\le a_r$ and $P\le1$. Lower bound: \cite[(3.6)]{Ba4} in the form of Step~7 above gives $X^{-1}=C^{(r)}=V^*G_{r+1}V-(\beta/a_r)P+a_r^{-1}\le\|V\|^2\|G_{r+1}\|+a_r^{-1}$ with $\|V\|\le1$ and $\|G_{r+1}\|\le L^2\gamma_0^{-1}$ (\Cref{rot:hyb:cor:L2}, the $(r+1)$-pin). Decay of $X$: \Cref{rot:hyb:prop:CT} for $A^X$ and the finite range of $P$ (its kernel vanishes unless $y,y'$ lie in one $L$-block, of diameter $<\sqrt3L$). Compression to $\Lambda$ keeps both bounds and the entries. Then, for a $1$-Lipschitz $\psi$ on $\Lambda$, $e^{t\psi}X|_\Lambda e^{-t\psi}-X|_\Lambda$ has entries bounded by $Ct\,d(y,y')e^{-(\theta/2)d(y,y')}$ for $t\le\theta/2$, hence norm $\le Ct$ (Schur test), so $e^{t\psi}X|_\Lambda e^{-t\psi}$ is invertible with inverse of norm $\le2(L^2\gamma_0^{-1}+a_r^{-1})$ for $t\le t_*$, $t_*$ depending on $a,L$ only; with $\psi=\min\{d(\cdot,y'),d(y,y')\}$ this gives $|\hat C^{(r)}_\Lambda(y,y')|\le Ce^{-t_*d(y,y')}$ (the entry is bounded by the operator norm on the counting measure of $\Gamma_r$). The difference with $C^{(r)}$ is $\hat C^{(r)}_\Lambda X_{\Lambda\Lambda^c}C^{(r)}_{\Lambda^c\Lambda}$ (the block identity for the inverse of a compression), whose decay follows from the three factors' decay and the geometric mean as in Step~6; the two-domain bound is the resolvent identity on $\ell^2(\Lambda)$ with \eqref{rot:hyb:eq:B2} for $\Delta^{(r)}$ ($\delta P=0$ off the collar), Step~5 and Step~6's interpolation.
\end{proof}
